\begin{abstractzh}

% TODO: 請在此處撰寫中文摘要（約 300-500 字）

本研究針對台灣高溫夏濕都市環境中行人熱壓力問題，建構一套整合物理資訊時空圖神經網路（PI-ST-GNN）、微氣候加權 $A^*$ 路徑演算法與多目標遺傳演算法（NSGA-II）的生成式都市設計最佳化框架。透過異質圖神經網路取代傳統高運算成本的計算流體力學（CFD）模擬，實現通用熱氣候指數（UTCI）之低延遲即時預測（驗證 $R^2 \approx 0.9966$），並以此驅動考量熱暴露成本的行人路徑規劃與 NSGA-II 多目標形態最佳化，同步輸出最佳化建築形態配置與涼適步行路徑，為氣候調適型都市設計提供即時決策支援。

關鍵詞：圖神經網路、都市微氣候、通用熱氣候指數、路徑最佳化、多目標遺傳演算法、生成式設計

\end{abstractzh}
